% AI-written, AI cross-reviewed; no human sign-off; not a claim that RH is proved
% RH Notes blueprint source excerpts for N076
% Verbatim excerpts in the notation of the parent blueprint. The labelled editions and source ranges appear below.
% See the unit YAML for the current proof status. Earlier-round passages retain their original wording.
\section*{N076. Bessel period in the wave grid}

% Parent source SHA-256: 44e3b079775c8d7cc0d43b0314bf43a3959be56b2bbafa7f26e5441a1a9163de
% blueprint/ch_T.tex lines 352 through 374; r6:N076
\subsection{N076. Removing the compact cutoff and separating equal types}\label{r6:N076}
\textit{Partial.} The threshold and the final weighted estimate are in Section~\ref{sec:T-r5-N076}. The pointwise normal form uses a positive local discriminant $G$. Its proof has two different limits, removal of a spatial cutoff and domination as the discriminant cutoff tends to zero. The two shifted planes enter only in the equal-type case.\footnote{\nolinkurl{thmBpAt_top}, \nolinkurl{thmBpAt_eqR_top} in \nolinkurl{CompanionTheory/AutomorphicTheory/CompactTopologicalForms/NormalFormTopDomination.lean}\allowbreak:10--35,50--76. \nolinkurl{nfTopDomAll_of_ptwC} in \nolinkurl{CompanionTheory/AutomorphicTheory/DiscreteGroupWaves/NFDomC.lean}\allowbreak:12--46,63--82.}

\paragraph{Extracting one coefficient.}
Fix an admissible cut word $w$ and a coefficient index in the word. Let $o$ bound the orders in its differential operator. In the model let $q=(q_1,\ldots,q_r,q_*)$, $d(q)=q_*\prod q_i$, and $G_q>0$. For exponent $s$ the evaluation of the word on $G_q^s f$ expands as $G_q^s d(q)^{-|w|}$ times a finite sum of evaluations on $f$. Each summand is a bounded multiplier followed by exactly $|w|$ operators in the model's hat set. Squaring the sum costs at most its number of terms by Cauchy--Schwarz. To extract the coefficient vectors $v_j$ of the differential operator, apply the Leibniz expansion to $G^s F$ for $s=0,\ldots,o$. If $V_s$ denotes its evaluated vector, the equations are
\[
 \sum_{j=0}^{o}s^j v_j=G_q^{-s}V_s .
\]
Terms of degree $>o$ vanish. The inverse Vandermonde matrix $(\lambda_{js})$ exists because $0,\ldots,o$ are distinct. Thus each squared coefficient norm is bounded by $\sum_{j,s}|\lambda_{js}|^2$ times $\sum_sG_q^{-2s}\|V_s\|^2$. The $G_q^{2s}$ from the squared model evaluation cancels this loss. All the resulting model operators and constants are selected before $F$, under the source's decay threshold.\footnote{\nolinkurl{model_pt}, \nolinkurl{normal_form_model_point} in \nolinkurl{CompanionTheory/AutomorphicTheory/JacobiDiagonalBounds/NFFinal.lean}\allowbreak:48--79. \nolinkurl{NFTieGen}, \nolinkurl{nfTieGen_of} in \nolinkurl{CompanionTheory/AutomorphicTheory/JacobiDiagonalBounds/NfTieHatSOf.lean}\allowbreak:63--85,87--115,182--200,212--243.}

\paragraph{Removing the cutoff.}
For each word the compact pointwise estimate uses finitely many differential operators $\Theta_j$ and a nonnegative multiplier $M$. The truncation bound for each $\Theta_j$ supplies an integrable $h_j$ dominating its estimates uniformly in the radial truncation $n$. Choose the common full-measure set for all $n,j$. At a local point with $G>0$, put $\eta=G/2$ and apply the compact estimate to $F_n=\mathrm{cut}_n\,\mathrm{cut}_{\eta}F$. The discriminant cutoff equals one on the open set $G>\eta$; consequently the differentiated coefficient there agrees with that of $\mathrm{cut}_n F$. Commuting the scalar cutoffs also identifies the top term with the radial approximation to the desired top term at the fixed discriminant cutoff $\varepsilon>0$. Pass to the pointwise limit in $n$. Outside the positive local locus the top term is zero. The resulting bound $M\sum_j |h_j|$ is integrable and independent of $\varepsilon$; measurability is supplied separately. The unequal-type endpoint uses this with the identity map and derivative order $2p_0$. The shifted endpoint invokes it on each of two planes with its own wave grid, lift grid and limiting bound. It requires order $2p_0+n_0$ on the original coefficient and reduces the decay budget from $A_0-(2r+2)p_0$ by another $n_0(r+1)$.\footnote{\nolinkurl{nfTopDomAll_of_ptwC} in \nolinkurl{CompanionTheory/AutomorphicTheory/DiscreteGroupWaves/NFDomC.lean}\allowbreak:12--52,53--82. \nolinkurl{thmBpAt_dom}, \nolinkurl{thmBpAt_eqR_dom} in \nolinkurl{CompanionTheory/AutomorphicTheory/CompactTopologicalForms/FinalC.lean}\allowbreak:21--52,67--103. \nolinkurl{thmBpAt_eqR_top} in \nolinkurl{CompanionTheory/AutomorphicTheory/CompactTopologicalForms/NormalFormTopDomination.lean}\allowbreak:55--76.}

\paragraph{Why two shifts suffice.}
Write the regular type in common-offset integer coordinates $L$, with $L_i-L_{i+1}\ge2n_0+1$ at the chosen gap and $L_0-L_{\rm last}\le2b$. In a branching summand with $L_{j+1}<M_j\le L_j$, split on $M_i\ge L_{i+1}+n_0+1$. In this case raise $L_{i+1}$ successively; otherwise lower $L_i$. At every intermediate step $j<n_0$ interlacing persists in the chosen direction. The one-step coefficient bound for that direction is $(3b)^{-2}$, using $j\le n_0\le b$ and the span bound. The product of the squared coefficients on the surviving chain is at least $(3b)^{-2n_0}$. The better chain can depend on $M$. On the spanning branching spaces the two operators $A_+A_+^*$ and $A_-A_-^*$ act by their chain scalars. The self-adjoint sum therefore has the same lower bound on all vectors. Equivariance and homogeneity spread the inequality from the last frame vector to an admissible $\omega$ orthogonal to the base vector:
\[
 (\operatorname{Re}H(\omega,\omega))^{n_0}\|v\|^2
 \le (3b)^{2n_0}\bigl(\|A_+(\omega)^*v\|^2+\|A_-(\omega)^*v\|^2\bigr).
\]
Use the corresponding shifts of $u'_{l+1}$ upward and $u'_l$ downward, selecting indices $l+1$ and $l$. Their transported coefficients are $f(g)A_+(g)$ and $f(g)A_-(g)$.\footnote{\nolinkurl{chain_prod_le}, \nolinkurl{chain_lower_bound} in \nolinkurl{Development/AbelianVarietyTheory/Proofs/JetHarmonics/ChainLowerBound.lean}\allowbreak:14--25,35--73,74--103. \nolinkurl{exists_shifted_types_lower_bound} in \nolinkurl{Development/AbelianVarietyTheory/Proofs/JetHarmonics/ExistsShiftedTypesLowerBound.lean}\allowbreak:19--65. \nolinkurl{shiftAtR_full} in \nolinkurl{CompanionTheory/AutomorphicTheory/BesselWaveGrids/EqShift.lean}\allowbreak:122--164.}

\paragraph{Outstanding propositions.}
\textit{[Proved in Lean, not reproduced.]} For the compact analytic step the required model identity is that the operator of \emph{each} admissible word and coefficient lies in the normal-form space with denominator operator $(\texttt{L.Di})^{|w|}$ and order $|w|$, with bounded multipliers commuting with the model functions. The translation to Bessel coefficients also uses the compatible gauge and the bound for products of hat operators on Jacobi transforms. On the representation side the unresolved one-step input identifies the adjoint norm on each branching component of the raised or lowered Pieri summand with its prescribed Wigner coefficient. The chain argument above uses this identification before multiplying the inequalities. \textit{[Hypothesis, discharged at the application.]} Ordinary and dual Zuckerman bounds and both clauses of \nolinkurl{WEnd.In91} are the original inputs from Section~\ref{sec:T-r5-N076}.\footnote{\nolinkurl{model_nf} in \nolinkurl{CompanionTheory/AutomorphicTheory/JacobiDiagonalBounds/NFFinal.lean}\allowbreak:19--32. \nolinkurl{normal_form_hat_compact} in \nolinkurl{CompanionTheory/AutomorphicTheory/JacobiDiagonalBounds/NormalFormHatCompact.lean}\allowbreak:7--14. \nolinkurl{normal_form_hat_product_bessel} in \nolinkurl{CompanionTheory/AutomorphicTheory/LocalBesselUniqueness/NFHatProdKBS.lean}\allowbreak:11--14. \nolinkurl{wigner_step_raise} in \nolinkurl{Development/AbelianVarietyTheory/Proofs/JetHarmonics/WignerStepRaisingUnitaryComplement.lean}\allowbreak:32--40. \nolinkurl{wigner_step_lowering} in \nolinkurl{Development/AbelianVarietyTheory/Proofs/JetHarmonics/WignerStepLoweringUnitaryComplement.lean}\allowbreak:35--43. \nolinkurl{wave_bessel_period_threshold} in \nolinkurl{CompanionTheory/AutomorphicTheory/BesselWaveGrids/WeightedBoundRotatedUniformOfCubicThreshold.lean}\allowbreak:23--28.}

